\section{Metodología}
\label{sec:metodologia}

\subsection{Diagramas}

Los diagramas exportados desde draw.io conviene guardarlos en PDF (vectorial) dentro de
\texttt{imagenes/}. También pueden dibujarse directamente con TikZ, como en la
\cref{fig:proceso}.

\begin{figure}[htbp]
  \centering
  \begin{tikzpicture}[node distance=6mm and 8mm,
      etapa/.style={draw=acento, fill=acento!8, rounded corners, minimum height=9mm,
                    text width=21mm, align=center, font=\small},
      flecha/.style={-Stealth, thick, acento}]
    \node[etapa] (sel) {Selección};
    \node[etapa, right=of sel] (pre) {Preprocesa\-miento};
    \node[etapa, right=of pre] (tra) {Transforma\-ción};
    \node[etapa, right=of tra] (min) {Minería de datos};
    \node[etapa, right=of min] (eva) {Interpreta\-ción};
    \foreach \a/\b in {sel/pre, pre/tra, tra/min, min/eva} \draw[flecha] (\a) -- (\b);
  \end{tikzpicture}
  \caption{Etapas del proceso \acs{KDD}.}
  \label{fig:proceso}
\end{figure}

Para incluir imágenes se usa \verb|\includegraphics| (la ruta \texttt{imagenes/} ya está
configurada). Con \texttt{subfigure} se agrupan varias en una misma figura
(\cref{fig:logos}).

\begin{figure}[htbp]
  \centering
  \begin{subfigure}[b]{0.3\textwidth}
    \centering
    \includegraphics[height=25mm]{utn}
    \caption{Logo a 25\,mm}
    \label{fig:logo-a}
  \end{subfigure}
  \hfill
  \begin{subfigure}[b]{0.3\textwidth}
    \centering
    \includegraphics[height=18mm]{utn}
    \caption{Logo a 18\,mm}
    \label{fig:logo-b}
  \end{subfigure}
  \caption{Ejemplo de subfiguras.}
  \label{fig:logos}
\end{figure}

\subsection{Código}

El código fuente se incluye desde archivos de la carpeta \texttt{codigo/}, de modo que el
informe siempre muestre la versión vigente (\cref{lst:sql,lst:python}).

\lstinputlisting[style=sql, caption={Ventas mensuales por categoría con funciones de ventana.},
                 label=lst:sql]{codigo/consulta.sql}

\lstinputlisting[style=python, caption={Validación cruzada de un árbol de decisión.},
                 label=lst:python]{codigo/modelo.py}

Fragmentos breves pueden escribirse en línea, como \lstinline[style=sql]|SELECT COUNT(*) FROM venta|,
o en un bloque:

\begin{lstlisting}[style=consola]
$ docker compose exec mysql mysql -u emmd_databases -p
\end{lstlisting}

\subsection{Algoritmos}

\begin{algorithm}[htbp]
  \caption{$k$-medias}
  \label{alg:kmedias}
  \KwIn{Conjunto de datos $\{\vect{x}_1, \dots, \vect{x}_n\}$, número de grupos $k$}
  \KwOut{Centroides $\vect{\mu}_1, \dots, \vect{\mu}_k$}
  Elegir $k$ centroides iniciales al azar\;
  \Repeat{los centroides no cambien}{
    \For{$i \gets 1$ \KwTo $n$}{
      $c_i \gets \argmin_{j} \norm{\vect{x}_i - \vect{\mu}_j}^2$\;
    }
    \For{$j \gets 1$ \KwTo $k$}{
      $\vect{\mu}_j \gets$ promedio de los $\vect{x}_i$ con $c_i = j$\;
    }
  }
\end{algorithm}

El \cref{alg:kmedias} se ejecutó para $k = 1, \dots, 8$.
